% Л14 -- EF Core: модель, миграции, запросы. Slides.
%
% Companion to the spoken script in ../L14-ef-core.md, which is the source of
% truth: the slides carry diagrams, code, real terminal output and short
% marks, and the delivery lives in the script only -- this deck carries no
% speaker notes. Slide text is in Russian because that is the content;
% everything else in this file is English.
%
% Build:  make          (latexmk -pdf; engine: pdflatex)
%
% Every C# listing was compiled on dotnet SDK 10.0 in a scratch
% `dotnet new console -n HelloWorld` with EF Core 10 and the Npgsql provider;
% SQL and tool output on the slides is captured from runs of that project
% against PostgreSQL 17. See ../CLAUDE.md for the verification table.
%
% House rules: ../CLAUDE.md and ../../../CLAUDE.md. Theme: ../../../../theme/.

\documentclass[aspectratio=169,10pt,t]{beamer}

\usetheme{dotnet}
\usetikzlibrary{arrows.meta,positioning,calc,decorations.pathreplacing,fit}
\usepackage{array}   % ragged-right table columns

% A ragged-right paragraph column: justified cells in narrow columns hyphenate
% badly, and a split word is the one thing the back row cannot read.
\newcolumntype{L}[1]{>{\raggedright\arraybackslash}p{#1}}

\dotnetlecture{Л14}

\title[Л14. EF Core: модель, миграции, запросы]{EF Core: модель, миграции, запросы}
\author{Дмитрий Курочкин}
\date{}

% Code sizes, one step down from \footnotesize. Smaller than 6.6pt does not
% carry to the back of the room.
\newcommand{\codesm}{\ttfamily\fontsize{7.6}{9.0}\selectfont}
\newcommand{\codexs}{\ttfamily\fontsize{7.0}{8.3}\selectfont}
\newcommand{\codexxs}{\ttfamily\fontsize{6.8}{7.8}\selectfont}
\newcommand{\termsm}{\ttfamily\fontsize{6.6}{7.9}\selectfont}

% A mark at the start of a listing line, for labels set beside the listing.
% Used through the listing's escapechar; the labels are drawn by an overlay
% picture after the listing, so the build needs two passes (latexmk does that).
\newcommand{\mk}[1]{\tikz[remember picture,overlay]\coordinate (#1) at (0,0.55ex);}

\tikzset{
  rawbrace/.style={decorate,decoration={brace,amplitude=4pt},hairline,
                   line width=0.8pt},
  rawlab/.style={anchor=west,font=\footnotesize,inner sep=0pt},
  box/.style={draw=hairline,line width=0.8pt,fill=codebg,inner sep=4pt,
              align=center,minimum height=0.9cm},
  arr/.style={-{Straight Barb[length=1.6mm]},muted,line width=0.7pt},
  arra/.style={-{Straight Barb[length=1.8mm]},accent,line width=0.9pt},
  arrl/.style={-{Straight Barb[length=1.8mm]},alt,line width=0.9pt},
  tbl/.style={draw=hairline,line width=0.8pt,fill=codebg,inner sep=3pt,
              align=left,font=\scriptsize\ttfamily,minimum width=2.2cm},
}

\begin{document}
\shorthandoff{"}

% ------------------------------------------------------------------ title --
% 1
\begin{frame}[plain]
  \titlepagednc{Лекция 14}{EF Core: модель, миграции, запросы}{}{Дмитрий Курочкин}
\end{frame}

% ============================================ 1. WHERE EF CORE BELONGS ==
\section{1. Место EF Core в приложении}

% 2
\begin{frame}[plain]
  \sectionopener{1}{Место EF Core в приложении}{Генератор кода, который вы
    \hi{обязаны проверять}.}
\end{frame}

% 3
\begin{frame}[plain]
  \vote{Что EF Core берёт на себя?}
\end{frame}

% 4
\begin{frame}{Чем EF Core не является}
  \slidetop
  \claim{От неё ждут трёх вещей, которых она \lo{не делает}.}

  {\renewcommand{\arraystretch}{1.6}
  \begin{tabular}{@{}L{4.4cm}L{8.4cm}@{}}
    \lo{\textbf{не кэш}}                 & загруженный объект живёт в контексте
      один запрос \\
    \lo{\textbf{не проектирование базы}} & индексы, ограничения, нормализация ---
      ваши решения \\
    \lo{\textbf{не замена SQL}}          & SQL пишете не вы, а проверяете --- вы \\
  \end{tabular}}
\end{frame}

% 5
\begin{frame}{Что берёт и что требует}
  \slidetop
  \versus{Берёт на себя}{%
    \begin{cbul}{accent}
      \item маппинг строк в объекты и обратно
      \item SQL из LINQ
      \item отслеживание изменений
      \item миграции схемы
    \end{cbul}}{Требует взамен}{%
    \begin{cbul}{alt}
      \item понимать SQL, который получился
      \item не считать базу деталью реализации
    \end{cbul}}

  \slidegap
  \hrulethin{0.42}
  \rulegap
  {\large ORM пишет SQL \hi{за вас}, но не \lo{вместо вас}.}
\end{frame}

% 6
\begin{frame}{Где доверять, где проверять}
  \slidetop
  \versus{Доверять}{%
    \begin{cbul}{accent}
      \item создание
      \item чтение по ключу
      \item обновление одной сущности
      \item удаление
    \end{cbul}}{Проверять каждый раз}{%
    \begin{cbul}{alt}
      \item запросы со связями
      \item группировки
      \item проекции с вычислениями
      \item больше одного оператора после \texttt{Where}
    \end{cbul}}

  \slidegap
  \hrulethin{0.42}
  \rulegap
  \aside{Ядро ORM общее, провайдер переводит в диалект: для PostgreSQL --- поверх
    драйвера Npgsql.}
\end{frame}

% ================================================ 2. MODEL AND CONFIG ==
\section{2. Модель и её конфигурация}

% 7
\begin{frame}[plain]
  \sectionopener{2}{Модель и её конфигурация}{Сущность --- обычный класс; всё
    про базу --- в \hi{отдельном классе} конфигурации.}
\end{frame}

% 8
\begin{frame}{Три слова}
  \slidetop
  {\renewcommand{\arraystretch}{1.6}
  \begin{tabular}{@{}L{4.0cm}L{8.8cm}@{}}
    \hi{\textbf{сущность}}        & класс, отображённый на таблицу; идентичность
      по ключу \\
    \hi{\textbf{DTO}}             & то, что уходит клиенту: без идентичности и без
      правил \\
    \hi{\textbf{доменная модель}} & сущности и типы-значения с правилами \\
  \end{tabular}}

  \slidegap
  \aside{Сущность ORM и доменная сущность в курсе --- \hi{один класс}; в больших
    системах их разделяют. DTO --- всегда отдельный тип, и наружу уходит только он.}
\end{frame}

% 9
\begin{frame}{Атрибуты или класс}
  \slidetop
  \versus{Класс конфигурации}{%
    \begin{cbul}{accent}
      \item домен не знает о хранении
      \item один класс --- одна таблица
      \item лежит в проекте данных
      \item подхватывается из сборки сам
    \end{cbul}}{Атрибуты на сущности}{%
    \begin{cbul}{alt}
      \item домен ссылается на пакет ORM
      \item мешают второму хранилищу
      \item тест домена тянет ORM
      \item цена видна не сразу
    \end{cbul}}

  \slidegap
  \hrulethin{0.42}
  \rulegap
  \aside{Цена приходит со вторым хранилищем и с тестом домена без ORM. Правило
    принимают \hi{заранее}.}
\end{frame}

% 10
\begin{frame}[fragile]{Сущность без атрибутов}
  \vspace{0.3em}
\begin{lstlisting}[style=csharp,basicstyle=\codexs,escapechar=|]
public sealed class Order
{
|\mk{e1}|    public int Id { get; private set; }
    public Email CustomerEmail { get; private set; } = null!;
    public decimal Total { get; private set; }
    public OrderStatus Status { get; private set; }
|\mk{e2}|    public List<OrderItem> Items { get; private set; } = [];

|\mk{e3}|    private Order() { }
|\mk{e4}|    public Order(Email customerEmail) =>
        CustomerEmail = customerEmail;

|\mk{e5}|    public void MarkPaid() => Status = Status == OrderStatus.New
        ? OrderStatus.Paid
        : throw new InvalidTransitionException(Status);
}
\end{lstlisting}
  \begin{tikzpicture}[remember picture,overlay,every node/.style={rawlab}]
    \node at ([xshift=10.4cm]e1) {запись закрыта};
    \node at ([xshift=10.4cm]e2) {позиции заказа};
    \node at ([xshift=10.4cm]e3) {для ORM: чтение из базы};
    \node at ([xshift=10.4cm]e4) {для домена};
    \node[text=accent] at ([xshift=10.4cm]e5) {правило перехода};
  \end{tikzpicture}

  \rulegap
  \aside{Ни одного слова про базу: домен \hi{не ссылается} на пакет ORM.}
\end{frame}

% 11
\begin{frame}[fragile]{Конфигурация --- отдельный класс}
  \vspace{0.3em}
\begin{lstlisting}[style=csharp,basicstyle=\codexs,escapechar=|]
public sealed class OrderConfiguration : IEntityTypeConfiguration<Order>
{
    public void Configure(EntityTypeBuilder<Order> builder)
    {
|\mk{k1}|        builder.ToTable("orders");
        builder.HasKey(o => o.Id);
|\mk{k2}|        builder.Property(o => o.Total).HasPrecision(18, 2);
|\mk{k3}|        builder.Property(o => o.Status)
            .HasConversion<string>().HasMaxLength(32);
|\mk{k4}|        builder.Property(o => o.CustomerEmail)
            .HasConversion(e => e.Value, v => new Email(v))
            .HasMaxLength(320);
|\mk{k5}|        builder.HasIndex(o => o.CustomerEmail);
|\mk{k6}|        builder.HasMany(o => o.Items).WithOne()
            .HasForeignKey(i => i.OrderId)
            .OnDelete(DeleteBehavior.Cascade);
    }
}
\end{lstlisting}
  \begin{tikzpicture}[remember picture,overlay,every node/.style={rawlab}]
    \node at ([xshift=10.6cm]k1) {таблица и ключ};
    \node at ([xshift=10.6cm]k2) {точность суммы};
    \node at ([xshift=10.6cm]k3) {статус строкой};
    \node[text=accent] at ([xshift=10.6cm]k4) {конвертер};
    \node at ([xshift=10.6cm]k5) {индекс};
    \node at ([xshift=10.6cm]k6) {один ко многим};
  \end{tikzpicture}
\end{frame}

% 12
\begin{frame}{Соглашения}
  \slidetop
  \versus{Задаём явно}{%
    \begin{cbul}{accent}
      \item ключ
      \item связь и правило удаления
      \item обязательность
      \item точность суммы
      \item длина строки, индекс
    \end{cbul}}{Соглашение сделает молча}{%
    \begin{cbul}{alt}
      \item \texttt{Id} --- ключ
      \item \texttt{OrderId} --- внешний ключ
      \item тип без \texttt{?} --- \texttt{NOT NULL}
      \item \texttt{decimal} --- точность базы
      \item \texttt{string} --- без длины
    \end{cbul}}

  \slidegap
  \hrulethin{0.42}
  \rulegap
  \aside{Соглашения удобны и опасны тем же: они \lo{невидимы}. Конфигурация
    читается как описание таблицы.}
\end{frame}

% 13
\begin{frame}[fragile]{Теневые свойства и глобальный фильтр}
  \slidetop
  \claim{Правило, которое нельзя обойти \hi{забывчивостью}.}

\begin{lstlisting}[style=csharp,basicstyle=\codesm,escapechar=|]
|\mk{s1}|builder.Property<DateTime>("UpdatedAt")
    .HasDefaultValueSql("now()");

|\mk{s2}|builder.HasQueryFilter(o => !o.IsDeleted);
\end{lstlisting}
  \begin{tikzpicture}[remember picture,overlay,every node/.style={rawlab}]
    \node at ([xshift=8.6cm]s1) {есть в таблице, нет в классе};
    \node[text=accent] at ([xshift=8.6cm]s2) {условие в каждом запросе};
  \end{tikzpicture}

  \slidegap
  \begin{cbul}{accent}
    \item теневое свойство заполняет перехватчик сохранения --- один на контекст
    \item так же ведут автора изменения и мягкое удаление
    \item фильтр отключается на одном запросе явно: \texttt{IgnoreQueryFilters()}
  \end{cbul}
\end{frame}

% 14
\begin{frame}[fragile]{Владеемый тип}
  \slidetop
  \claim{Одно свойство в коде --- \hi{две колонки} той же таблицы.}

\begin{lstlisting}[style=csharp,basicstyle=\codesm,escapechar=|]
|\mk{w1}|public record Money(decimal Amount, string Currency);

|\mk{w2}|builder.OwnsOne(o => o.Delivery, money =>
{
|\mk{w3}|    money.Property(m => m.Amount).HasPrecision(18, 2);
|\mk{w4}|    money.Property(m => m.Currency).HasMaxLength(3);
});
\end{lstlisting}
  \begin{tikzpicture}[remember picture,overlay,every node/.style={rawlab}]
    \node at ([xshift=9.4cm]w1) {доменный тип из двух полей};
    \node[text=accent] at ([xshift=9.4cm]w2) {живёт внутри сущности};
    \node at ([xshift=9.4cm]w3) {\texttt{delivery\_amount}};
    \node at ([xshift=9.4cm]w4) {\texttt{delivery\_currency}};
  \end{tikzpicture}

  \slidegap
  \aside{Ни отдельной таблицы, ни ключа. Конвертер --- для типа из одного поля,
    владеемый тип --- из нескольких.}
\end{frame}

% 15
\begin{frame}[fragile]{Имена}
  \slidetop
  \claim{PostgreSQL приводит имена без кавычек к \hi{нижнему регистру}.}

  \eyebrow{по умолчанию}
\begin{lstlisting}[style=term,basicstyle=\termsm]
SELECT o."Id", o."CustomerEmail", o."Status", o."Total"
FROM orders AS o
WHERE o."CustomerEmail" = @email
\end{lstlisting}

  \rulegap
  \eyebrow{Use Snake Case Naming Convention}
\begin{lstlisting}[style=term,basicstyle=\termsm]
SELECT o.id, o.customer_email, o.status, o.total
FROM orders AS o
WHERE o.customer_email = @email
\end{lstlisting}

  \rulegap
  \aside{Одна строка на всю модель. Вручную имя задают только исключениям ---
    и только в конфигурации.}
\end{frame}

% 16
\begin{frame}{Как генерируется ключ}
  \slidetop
  {\renewcommand{\arraystretch}{1.5}
  \begin{tabular}{@{}L{3.6cm}L{4.7cm}L{4.5cm}@{}}
    & \colhead{accent}{даёт} & \colhead{alt}{цена} \\
    \textbf{автоинкремент}      & просто и компактно           & известен только
      после вставки \\
    \textbf{последовательность} & номер можно взять заранее    & ещё одна поездка
      к базе \\
    \textbf{на клиенте}         & известен до сохранения, не раскрывает количество
      & длиннее; случайный --- хуже для индекса \\
  \end{tabular}}

  \slidegap
  \aside{В курсе --- \hi{автоинкремент}. Суррогатный ключ из одной колонки на
    каждой таблице; естественная уникальность --- отдельный уникальный индекс.}
\end{frame}

% 17
\begin{frame}{Три вида связей}
  \slidetop
  \begin{center}
    \begin{tikzpicture}
      % 1:1
      \node[tbl] (a1) at (0,0) {orders\\\textbf{id}};
      \node[tbl] (a2) at (0,-1.7) {invoices\\\textbf{order\_id}};
      \draw[arra] (a2) -- (a1);
      \node[font=\footnotesize\bfseries,text=accent] at (0,1.0) {1 : 1};
      % 1:many
      \node[tbl] (b1) at (4.6,0) {orders\\\textbf{id}};
      \node[tbl] (b2) at (4.6,-1.7) {order\_items\\id\\\textbf{order\_id}};
      \draw[arra] (b2) -- (b1);
      \node[font=\footnotesize\bfseries,text=accent] at (4.6,1.0) {1 : many};
      % many:many
      \node[tbl] (c1) at (8.2,0) {orders\\\textbf{id}};
      \node[tbl] (c3) at (11.6,0) {tags\\\textbf{id}};
      \node[tbl] (c2) at (9.9,-1.7) {order\_tag\\\textbf{order\_id}\\\textbf{tag\_id}};
      \draw[arra] (c2) -- (c1);
      \draw[arra] (c2) -- (c3);
      \node[font=\footnotesize\bfseries,text=accent] at (9.9,1.0) {many : many};
    \end{tikzpicture}
  \end{center}

  \slidegap
  {\small
  \begin{tabular}{@{}L{4.0cm}L{3.9cm}L{4.6cm}@{}}
    ключ одной --- внешний ключ на другую & внешний ключ в таблице «многих» &
      промежуточная таблица с двумя внешними ключами \\
  \end{tabular}}
\end{frame}

% 18
\begin{frame}[fragile]{Один ко многим}
  \slidetop
  \claim{Правило удаления --- \hi{решение}, а не умолчание.}

\begin{lstlisting}[style=csharp,basicstyle=\codesm,escapechar=|]
|\mk{m1}|builder.HasMany(o => o.Items)
|\mk{m2}|    .WithOne()
|\mk{m3}|    .HasForeignKey(i => i.OrderId)
|\mk{m4}|    .OnDelete(DeleteBehavior.Cascade);
\end{lstlisting}
  \begin{tikzpicture}[remember picture,overlay,every node/.style={rawlab}]
    \node at ([xshift=7.6cm]m1) {коллекция у заказа};
    \node at ([xshift=7.6cm]m2) {обратной навигации нет};
    \node at ([xshift=7.6cm]m3) {\texttt{int} без \texttt{?} --- связь обязательна};
    \node[text=accent] at ([xshift=7.6cm]m4) {позиции без заказа не живут};
  \end{tikzpicture}

  \slidegap
  \begin{cbul}{alt}
    \item удалили покупателя --- каскадом ушли его заказы и \lo{история}
    \item навигация «на всякий случай» --- цикл для сериализатора
  \end{cbul}
\end{frame}

% 19
\begin{frame}[plain]
  \vote{Почта --- не строка, а тип с проверкой. В базе колонка почты ---
    текст. Что делать?}
\end{frame}

% 20
\begin{frame}[fragile]{Конвертер}
  \slidetop
  \claim{Свойство --- \hi{доменный тип}, колонка --- примитив.}

\begin{lstlisting}[style=csharp,basicstyle=\codesm,escapechar=|]
builder.Property(o => o.CustomerEmail)
|\mk{c1}|    .HasConversion(e => e.Value, v => new Email(v));

|\mk{c2}|context.Orders.Where(o => o.CustomerEmail == email)
\end{lstlisting}
  \begin{tikzpicture}[remember picture,overlay,every node/.style={rawlab}]
    \node[text=accent] at ([xshift=9.2cm]c1) {в базу и из базы};
    \node at ([xshift=9.2cm]c2) {сравнение по доменному типу};
  \end{tikzpicture}

\begin{lstlisting}[style=term,basicstyle=\termsm]
-- @email='anna@example.com' (Nullable = false)
SELECT o.id, o.customer_email, o.status, o.total
FROM orders AS o
WHERE o.customer_email = @email
\end{lstlisting}

  \rulegap
  \aside{Равенство переводится в SQL. Операция \lo{внутри} значения --- «заканчивается
    на домен» --- по доменному типу не переводится.}
\end{frame}

% 21
\begin{frame}[plain]
  \keyline{Сущность --- обычный класс без атрибутов.\\ Всё про базу --- в классе
    конфигурации, по одному на сущность.}
\end{frame}

% ========================================================= 3. DBCONTEXT ==
\section{3. DbContext}

% 22
\begin{frame}[plain]
  \sectionopener{3}{DbContext}{Сессия работы с базой на \hi{один запрос}, и
    ничего больше.}
\end{frame}

% 23
\begin{frame}[fragile]{Контекст целиком}
  \vspace{0.3em}
\begin{lstlisting}[style=csharp,basicstyle=\codexs,escapechar=|]
public sealed class OrdersDbContext(DbContextOptions<OrdersDbContext> options)
    : DbContext(options)
{
|\mk{x1}|    public DbSet<Order> Orders => Set<Order>();
    public DbSet<Customer> Customers => Set<Customer>();

    protected override void OnModelCreating(ModelBuilder modelBuilder) =>
|\mk{x2}|        modelBuilder.ApplyConfigurationsFromAssembly(
            typeof(OrdersDbContext).Assembly);
}

|\mk{x3}|services.AddDbContext<OrdersDbContext>(options =>
|\mk{x4}|    options.UseNpgsql(configuration.GetConnectionString("Orders"))
|\mk{x5}|           .UseSnakeCaseNamingConvention());
\end{lstlisting}
  \begin{tikzpicture}[remember picture,overlay,every node/.style={rawlab}]
    \node at ([xshift=10.6cm]x1) {набор на таблицу};
    \node[text=accent] at ([xshift=10.6cm]x2) {все конфигурации разом};
    \node at ([xshift=10.6cm]x3) {scoped};
    \node at ([xshift=10.6cm]x4) {строка из конфигурации};
    \node at ([xshift=10.6cm]x5) {змеиный регистр};
  \end{tikzpicture}

  \slidegap
  \aside{Единица работы: всё изменённое в одном контексте сохраняется одной
    транзакцией, \hi{одним вызовом}.}
\end{frame}

% 24
\begin{frame}{Чего в контексте нет}
  \slidetop
  \claim{Наборы и подключение конфигураций. \hi{Всё.}}

  {\renewcommand{\arraystretch}{1.6}
  \begin{tabular}{@{}L{4.0cm}L{8.8cm}@{}}
    \lo{\textbf{бизнес-логики}} & можно ли отменить заказ, решает сущность или
      сервис \\
    \lo{\textbf{запросов}}      & в хранилище: методы названы по сценариям,
      параметры простые \\
    \lo{\textbf{конфигураций}}  & в отдельных классах, по одному на сущность \\
  \end{tabular}}

  \slidegap
  \aside{Тонкое хранилище прячет ORM \hi{от домена}, а не от вас. «Хранилище чего
    угодно» с лямбдой в параметре --- ORM через заднюю дверь.}
\end{frame}

% 25
\begin{frame}[plain]
  \vote{Контекст регистрируется как scoped. Почему не singleton --- он же
    дорогой?}
\end{frame}

% 26
\begin{frame}{Время жизни}
  \slidetop
  {\renewcommand{\arraystretch}{1.6}
  \begin{tabular}{@{}L{3.2cm}L{9.6cm}@{}}
    \lo{\textbf{singleton}} & одно соединение на всех; отслеживаемые объекты
      копятся вечно; не потокобезопасен \\
    \lo{\textbf{transient}} & два хранилища в запросе --- два контекста --- две
      транзакции \\
    \hi{\textbf{scoped}}    & один на запрос: одна единица работы, все хранилища
      получают его же \\
  \end{tabular}}

  \slidegap
  \aside{Контекст не дорогой: дорогая модель, и она строится один раз. Где
    запроса нет --- область на итерацию или фабрика контекстов.}
\end{frame}

% 27
\begin{frame}[fragile]{Раскладка проекта}
  \vspace{0.3em}
\begin{lstlisting}[style=plain,basicstyle=\codexs,escapechar=|]
OrdersService.Data/
|\mk{r1}|    OrdersDbContext.cs
|\mk{r2}|    Configurations/
        OrderConfiguration.cs
        OrderItemConfiguration.cs
        CustomerConfiguration.cs
|\mk{r3}|    Migrations/
        20260901120000_Initial.cs
        20260910093000_AddOrderStatusIndex.cs
        OrdersDbContextModelSnapshot.cs
|\mk{r4}|    Repositories/
        IOrderRepository.cs
        OrderRepository.cs
|\mk{r5}|    DataModule.cs
\end{lstlisting}
  \begin{tikzpicture}[remember picture,overlay,every node/.style={rawlab}]
    \node at ([xshift=7.6cm]r1) {контекст в корне};
    \node at ([xshift=7.6cm]r2) {по файлу на сущность};
    \node[text=accent] at ([xshift=7.6cm]r3) {генерирует инструмент; коммитятся};
    \node at ([xshift=7.6cm]r4) {интерфейс может жить в домене};
    \node at ([xshift=7.6cm]r5) {регистрация одной строкой};
  \end{tikzpicture}

  \slidegap
  \aside{Три проекта: домен, доступ к данным, веб. Ссылки --- \hi{в одну сторону}:
    домен не ссылается ни на что.}
\end{frame}

% 28
\begin{frame}[plain]
  \keyline{Контекст --- наборы плюс подключение конфигураций.\\ Scoped, один на
    запрос, и ничего больше.}
\end{frame}

% ======================================================== 4. MIGRATIONS ==
\section{4. Миграции}

% 29
\begin{frame}[plain]
  \sectionopener{4}{Миграции}{Схема как код: изменение базы живёт \hi{рядом} с
    изменением модели.}
\end{frame}

% 30
\begin{frame}{Как устроены миграции}
  \slidetop
  \begin{center}
    \begin{tikzpicture}[font=\footnotesize,
      tb/.style={box,minimum width=2.9cm,minimum height=1.1cm}]
      \node[tb] (m1) at (0,0)    {\textbf{модель}\\сущности и конфигурации};
      \node[tb] (m2) at (0,-1.6) {\textbf{снимок}\\модель в прошлый раз};
      \node[tb,draw=accent,line width=1.0pt] (m3) at (4.9,-0.8)
        {\textbf{класс миграции}\\\texttt{Up} и \texttt{Down}};
      \node[tb] (m4) at (9.8,-0.8) {\textbf{таблица истории}\\что уже применено};
      \draw[arr] (m1.east) -- (m3.west);
      \draw[arr] (m2.east) -- (m3.west);
      \draw[arra] (m3) -- node[above,font=\scriptsize,text=accent] {применить} (m4);
      \node[font=\scriptsize,text=muted] at (2.45,-0.8) {разница};
    \end{tikzpicture}
  \end{center}

  \slidegap
  \begin{cbul}{accent}
    \item класс и снимок коммитятся \hi{вместе} с изменением сущности
    \item применить --- выполнить все, которых нет в таблице истории, по порядку
    \item снимок руками не редактируют: конфликт --- пересоздать миграцию
  \end{cbul}
\end{frame}

% 31
\begin{frame}[fragile]{Рабочий цикл}
  \slidetop
\begin{lstlisting}[style=term,basicstyle=\termsm]
$ dotnet ef migrations add AddOrderStatusIndex
Build started...
Build succeeded.
Done. To undo this action, use 'ef migrations remove'
$ dotnet ef migrations list
Build started...
Build succeeded.
20261112060351_Initial
20261112060513_AddOrderStatusIndex (Pending)
\end{lstlisting}

\begin{lstlisting}[style=csharp,basicstyle=\codexs]
migrationBuilder.CreateIndex(
    name: "ix_orders_status", table: "orders", column: "status");
\end{lstlisting}

  \rulegap
  \aside{Изменили модель --- добавили --- \hi{прочитали} --- применили. Имя ---
    описание изменения. Применённую миграцию \lo{не редактируют}.}
\end{frame}

% 32
\begin{frame}[fragile]{Что здесь не так?}
  \slidetop
\begin{lstlisting}[style=csharp,basicstyle=\codesm]
public partial class AddOrderArchive : Migration
{
    protected override void Up(MigrationBuilder migrationBuilder)
    {
        migrationBuilder.AddColumn<string>(
            name: "metadata", table: "orders", type: "jsonb", nullable: true);

        migrationBuilder.Sql(
            "ALTER TABLE orders ADD COLUMN archived_at TIMESTAMPTZ; " +
            "CREATE INDEX CONCURRENTLY ix_orders_archived ON orders (archived_at);");
    }

    protected override void Down(MigrationBuilder migrationBuilder) { /* ... */ }
}
\end{lstlisting}
\end{frame}

% 33
\begin{frame}[fragile]{Три ошибки}
  \vspace{0.3em}
\begin{lstlisting}[style=csharp,basicstyle=\codexs]
migrationBuilder.AddColumn<string>(
    name: "metadata", table: "orders", type: "jsonb", nullable: true);
migrationBuilder.Sql(
    "ALTER TABLE orders ADD COLUMN archived_at TIMESTAMPTZ; " +
    "CREATE INDEX CONCURRENTLY ix_orders_archived ON orders (archived_at);");
\end{lstlisting}

\begin{lstlisting}[style=term,basicstyle=\termsm,breaklines=true,breakatwhitespace=true,breakindent=4em]
$ dotnet ef database update
...
Npgsql.PostgresException (0x80004005): 25001: CREATE INDEX CONCURRENTLY cannot run inside a transaction block
\end{lstlisting}

  \rulegap
  {\renewcommand{\arraystretch}{1.3}
  \begin{tabular}{@{}L{3.2cm}L{9.6cm}@{}}
    \lo{\textbf{тип строкой}} & \texttt{jsonb} --- тип одной базы: миграция
      привязана к провайдеру \\
    \lo{\textbf{сырой SQL}}   & диалект; колонка не попадёт в снимок --- модель
      разойдётся со схемой \\
    \lo{\textbf{транзакция}}  & миграция идёт в транзакции --- падает и на
      PostgreSQL \\
  \end{tabular}}
\end{frame}

% 34
\begin{frame}{Как правильно}
  \slidetop
  \versus{В конфигурации}{%
    \begin{cbul}{accent}
      \item тип колонки --- средствами провайдера
      \item индекс --- \texttt{HasIndex}
      \item миграция выводится из модели
      \item всё попадает в снимок
    \end{cbul}}{Сырой SQL}{%
    \begin{cbul}{alt}
      \item только то, чего построитель не умеет
      \item перенос данных, справочник
      \item по стандарту SQL
      \item диалект --- под проверкой провайдера
    \end{cbul}}

  \slidegap
  \hrulethin{0.42}
  \rulegap
  \aside{Миграция быстрая. Заполнить колонку на ста миллионах строк --- фоновая
    задача, порциями, после выкладки.}
\end{frame}

% 35
\begin{frame}[plain]
  \vote{При старте приложение само применяет миграции. Что не так?}
\end{frame}

% 36
\begin{frame}{Кто применяет миграции}
  \slidetop
  \versus{Шаг поставки}{%
    \begin{cbul}{accent}
      \item один раз, до выкладки
      \item из конвейера сборки
      \item SQL-скрипт или исполняемый файл
      \item роль с правом менять схему
    \end{cbul}}{Приложение на старте}{%
    \begin{cbul}{alt}
      \item два экземпляра --- две попытки
      \item вторая падает или удваивает данные
      \item приложению нужны права на схему
      \item ошибка в коде --- ошибка в схеме
    \end{cbul}}

  \slidegap
  \hrulethin{0.42}
  \rulegap
  \aside{Приложение стартует уже на новой схеме, с ролью, которая умеет
    \hi{только данные}.}
\end{frame}

% 37
\begin{frame}[fragile]{Идемпотентный скрипт}
  \slidetop
\begin{lstlisting}[style=term,basicstyle=\termsm,breaklines=true,breakatwhitespace=true,breakindent=4em]
$ dotnet ef migrations script Initial --idempotent
Build started...
Build succeeded.
START TRANSACTION;

DO $EF$
BEGIN
    IF NOT EXISTS(SELECT 1 FROM "__EFMigrationsHistory" WHERE "migration_id" = '20261112060513_AddOrderStatusIndex') THEN
    CREATE INDEX ix_orders_status ON orders (status);
    END IF;
END $EF$;
\end{lstlisting}

  \rulegap
  \begin{cbul}{accent}
    \item скрипт можно \hi{прочитать} и отдать администратору базы
    \item \texttt{migrations bundle} --- исполняемый файл: без исходников и SDK
    \item \texttt{has-pending-model-changes} --- забытая миграция не пройдёт
  \end{cbul}
\end{frame}

% 38
\begin{frame}{Выкладка без остановки}
  \slidetop
  \claim{Пока идёт выкладка, старый код живёт на \hi{новой схеме}.}

  {\renewcommand{\arraystretch}{1.5}
  \begin{tabular}{@{}L{4.4cm}L{8.4cm}@{}}
    \hi{\textbf{добавить колонку}}      & старый код не ломает \\
    \textbf{новая обязательная}         & три шага: допускает \texttt{null} ---
      заполнить --- обязательная \\
    \textbf{переименовать}              & четыре шага, две выкладки: добавить,
      копировать, переключить код, удалить \\
    \lo{\textbf{удалить колонку}}       & две выкладки: код перестал использовать
      --- миграция удалила \\
  \end{tabular}}

  \slidegap
  \aside{Инструмент сгенерирует переименование одной командой --- и об этом
    \lo{не знает}. Читайте миграции глазами.}
\end{frame}

% 39
\begin{frame}[plain]
  \keyline{Миграция выводится из модели и не знает провайдера.\\ Применяется
    отдельным шагом до выкладки, а не на старте.}
\end{frame}

% =========================================================== 5. QUERIES ==
\section{5. Запросы и загрузка данных}

% 40
\begin{frame}[plain]
  \sectionopener{5}{Запросы и загрузка данных}{Чтение --- девяносто процентов
    работы с базой и \hi{сто} процентов сюрпризов.}
\end{frame}

% 41
\begin{frame}{Три способа загрузить связи}
  \slidetop
  {\renewcommand{\arraystretch}{1.6}
  \begin{tabular}{@{}L{2.6cm}L{4.4cm}L{5.8cm}@{}}
    \hi{\textbf{eager}}  & вместе, одним запросом   & \texttt{Include} --- вы знаете
      заранее, что нужно \\
    \textbf{explicit}    & потом, отдельным вызовом & нужно ли --- решается по ходу \\
    \lo{\textbf{lazy}}   & сама, при обращении      & свойство идёт в базу:
      \lo{N+1} \\
  \end{tabular}}

  \slidegap
  \aside{\texttt{Include} умеет вложенность и фильтр: условие уходит в
    соединение, а не в память.}
\end{frame}

% 42
\begin{frame}[plain]
  \vote{Загрузили сто заказов и в цикле обратились к позициям каждого. Сколько
    запросов уйдёт в базу при ленивой загрузке?}
\end{frame}

% 43
\begin{frame}[fragile]{N+1}
  \slidetop
  \claim{Обращение к свойству выглядит как \lo{чтение поля}.}

\begin{lstlisting}[style=csharp,basicstyle=\codesm,escapechar=|]
|\mk{n1}|var orders = await context.Orders.ToListAsync(ct);

int items = 0;
foreach (var order in orders)
|\mk{n2}|    items += order.Items.Count;
\end{lstlisting}
  \begin{tikzpicture}[remember picture,overlay,every node/.style={rawlab}]
    \node at ([xshift=8.8cm]n1) {один запрос за заказами};
    \node[text=alt] at ([xshift=8.8cm]n2) {и по одному на каждый заказ};
  \end{tikzpicture}

  \rulegap
  {\renewcommand{\arraystretch}{1.3}
  \begin{tabular}{@{}L{3.4cm}l@{}}
    \lo{ленивая загрузка}       & {\codexs orders: 100, items: 200, commands: 101} \\
    \hi{\texttt{Include}}       & {\codexs orders: 100, items: 200, commands: 1} \\
  \end{tabular}}

  \rulegap
  \aside{На дашборде --- гребёнка одинаковых отрезков. Ленивую загрузку в сервисе
    \lo{не включать}.}
\end{frame}

% 44
\begin{frame}[fragile]{Две коллекции в Include}
  \slidetop
  \claim{Соединение двух коллекций --- \lo{декартово произведение}.}

\begin{lstlisting}[style=csharp,basicstyle=\codesm,escapechar=|]
context.Orders
|\mk{u1}|    .Include(o => o.Items)
|\mk{u2}|    .Include(o => o.Tags)
|\mk{u3}|    .AsSplitQuery()
\end{lstlisting}
  \begin{tikzpicture}[remember picture,overlay,every node/.style={rawlab}]
    \node at ([xshift=7.6cm]u1) {10 позиций};
    \node[text=alt] at ([xshift=7.6cm]u2) {и 5 тегов --- 50 строк на заказ};
    \node[text=accent] at ([xshift=7.6cm]u3) {запрос на коллекцию, без умножения};
  \end{tikzpicture}

  \slidegap
  \begin{cbul}{accent}
    \item тысяча позиций и сотня тегов --- сто тысяч строк ради тысячи ста
    \item одна коллекция в \texttt{Include} --- обычный запрос; несколько ---
      \hi{разделённый}
  \end{cbul}
\end{frame}

% 45
\begin{frame}[fragile]{Проекция вместо сущности}
  \vspace{0.3em}
\begin{lstlisting}[style=csharp,basicstyle=\codexs,escapechar=|]
|\mk{p1}|var orders = await context.Orders
|\mk{p2}|    .Include(o => o.Items)
    .Where(o => o.CustomerEmail == email)
|\mk{p3}|    .ToListAsync(ct);
var summaries = orders
|\mk{p4}|    .Select(o => new OrderSummaryDto(o.Id, o.Status, o.Items.Count))
    .ToList();

|\mk{p5}|var summaries = await context.Orders
    .Where(o => o.CustomerEmail == email)
|\mk{p6}|    .Select(o => new OrderSummaryDto(o.Id, o.Status, o.Items.Count))
    .ToListAsync(ct);
\end{lstlisting}
  \begin{tikzpicture}[remember picture,overlay,every node/.style={rawlab}]
    \node[text=alt] at ([xshift=10.8cm]p1) {\textbf{было}};
    \node at ([xshift=10.8cm]p2) {все позиции целиком};
    \node at ([xshift=10.8cm]p3) {всё отслеживается};
    \node at ([xshift=10.8cm]p4) {DTO --- в памяти};
    \node[text=accent] at ([xshift=10.8cm]p5) {\textbf{стало}};
    \node at ([xshift=10.8cm]p6) {DTO --- в запросе};
  \end{tikzpicture}

  \slidegap
  \aside{Читаете, чтобы показать, --- \hi{проецируйте}. Сущность целиком --- только
    чтобы её менять.}
\end{frame}

% 46
\begin{frame}[fragile]{Два SQL}
  \slidetop
  \eyebrow{было}
\begin{lstlisting}[style=term,basicstyle=\termsm]
SELECT o.id, o.customer_email, o.status, o.total, o0.id, o0.order_id, o0.quantity, o0.sku
FROM orders AS o
LEFT JOIN order_items AS o0 ON o.id = o0.order_id
WHERE o.customer_email = @email
ORDER BY o.id
\end{lstlisting}

  \rulegap
  \eyebrow{стало}
\begin{lstlisting}[style=term,basicstyle=\termsm]
SELECT o.id, o.status, (
    SELECT count(*)::int
    FROM order_items AS o0
    WHERE o.id = o0.order_id)
FROM orders AS o
WHERE o.customer_email = @email
\end{lstlisting}

  \rulegap
  \aside{Разница в SQL --- и есть разница в \hi{нагрузке на базу}.}
\end{frame}

% 47
\begin{frame}[fragile]{Читаем --- не отслеживаем}
  \slidetop
\begin{lstlisting}[style=csharp,basicstyle=\codesm,escapechar=|]
var paid = context.Orders
|\mk{a1}|    .AsNoTracking()
    .Where(o => o.Status == OrderStatus.Paid)
|\mk{a2}|    .OrderBy(o => o.Id)
|\mk{a3}|    .AsAsyncEnumerable();

|\mk{a4}|await foreach (var order in paid.WithCancellation(ct))
    Console.WriteLine(order.Id);
\end{lstlisting}
  \begin{tikzpicture}[remember picture,overlay,every node/.style={rawlab}]
    \node[text=accent] at ([xshift=9.2cm]a1) {без отслеживателя};
    \node at ([xshift=9.2cm]a2) {порядок однозначный};
    \node at ([xshift=9.2cm]a3) {по строке, а не списком};
    \node at ([xshift=9.2cm]a4) {токен доходит до базы};
  \end{tikzpicture}

  \slidegap
  \begin{cbul}{accent}
    \item для запросов на чтение \texttt{AsNoTracking} --- норма
    \item без отслеживания нет тождественности: одна строка --- \lo{два объекта}
    \item запрос без токена доработает и после ухода клиента
  \end{cbul}
\end{frame}

% 48
\begin{frame}[plain]
  \vote{У набора есть «найти по ключу». Чем он отличается от «первый, у кого
    идентификатор равен»?}
\end{frame}

% 49
\begin{frame}{Три способа увидеть SQL}
  \slidetop
  \claim{Написали запрос --- посмотрели SQL --- \hi{сравнили} с ожиданием.}

  {\renewcommand{\arraystretch}{1.6}
  \begin{tabular}{@{}L{3.6cm}L{9.2cm}@{}}
    \hi{\textbf{лог}}                    & категория команд базы на уровне
      \texttt{Information}: текст, параметры, длительность \\
    \hi{\textbf{\texttt{ToQueryString()}}} & текст SQL до выполнения: в отладчике и
      в тесте \\
    \hi{\textbf{дашборд}}                & отрезок на каждый запрос; видно, сколько
      их на один HTTP-запрос \\
  \end{tabular}}

  \slidegap
  \aside{Значения параметров в логе --- только в разработке: это персональные
    данные. SQL верный, а запрос медленный --- это \lo{план}, а не ORM.}
\end{frame}

% 50
\begin{frame}[plain]
  \keyline{Проецируйте в DTO в запросе, сущности загружайте для изменения.\\
    Смотрите SQL после каждого нового запроса.}
\end{frame}

% ============================================= 6. RULES AND ALTERNATIVES ==
\section{6. Правила, грабли и альтернативы}

% 51
\begin{frame}[plain]
  \sectionopener{6}{Правила, грабли и альтернативы}{Шесть тем о том, где
    ORM начинает \lo{мешать}.}
\end{frame}

% 52
\begin{frame}{Отслеживание изменений}
  \slidetop
  \claim{Состояние живёт \lo{между строчками} вашего кода.}

  \begin{cbul}{alt}
    \item метод изменил заказ и не сохранил --- сохранил другой метод, вместе со
      своим
    \item объект отдали наружу --- чужое изменение «на всякий случай» сохранилось
    \item загрузили без отслеживания, изменили --- не сохранилось
  \end{cbul}

  \slidegap
  \hrulethin{0.42}
  \rulegap
  {\large Изменение и сохранение --- \hi{в одном методе}, рядом.}

  \rulegap
  \aside{Меняют методом сущности с правилом, а не сеттером. Record --- для DTO:
    \texttt{with} создаёт новый объект, отслеживатель следит за старым.}
\end{frame}

% 53
\begin{frame}[fragile]{Транзакции и гонки}
  \slidetop
  \claim{Условие «где версия та же» --- тихая перезапись \hi{невозможна}.}

\begin{lstlisting}[style=csharp,basicstyle=\codesm]
builder.Property<uint>("Version").IsRowVersion();
\end{lstlisting}

\begin{lstlisting}[style=term,basicstyle=\termsm]
      UPDATE orders SET status = @p0
      WHERE id = @p1 AND xmin = @p2
      RETURNING xmin;
\end{lstlisting}

  \rulegap
  \begin{cbul}{accent}
    \item условие не сработало --- \texttt{DbUpdateConcurrencyException}, наружу
      \hi{409}
    \item сохранение --- одна транзакция: проходит целиком или не проходит
    \item несколько сохранений --- транзакция явно, на контексте
    \item уровень изоляции --- по сценарию, а не «самый строгий на всякий случай»
  \end{cbul}
\end{frame}

% 54
\begin{frame}[fragile]{Массовые операции}
  \vspace{0.3em}
\begin{lstlisting}[style=csharp,basicstyle=\codexxs,escapechar=|]
|\mk{z1}|var stale = await context.Orders
    .Where(o => o.Status == OrderStatus.New && o.CreatedAt < cutoff)
    .ToListAsync(ct);
foreach (var order in stale) order.Cancel();
|\mk{z2}|await context.SaveChangesAsync(ct);

|\mk{z3}|await context.Orders
    .Where(o => o.Status == OrderStatus.New && o.CreatedAt < cutoff)
    .ExecuteUpdateAsync(
        s => s.SetProperty(o => o.Status, OrderStatus.Cancelled), ct);

|\mk{z4}|options.UseNpgsql(cs, npgsql => npgsql.EnableRetryOnFailure());
\end{lstlisting}
  \begin{tikzpicture}[remember picture,overlay,every node/.style={rawlab}]
    \node[text=alt] at ([xshift=10.2cm]z1) {\textbf{было}: N объектов};
    \node at ([xshift=10.2cm]z2) {и N команд};
    \node[text=accent] at ([xshift=10.2cm]z3) {\textbf{стало}: одна команда};
    \node at ([xshift=10.2cm]z4) {повтор при сбое базы};
  \end{tikzpicture}

\begin{lstlisting}[style=term,basicstyle=\termsm]
      UPDATE orders AS o
      SET status = @p
      WHERE o.status = 'New' AND o.created_at < @cutoff
\end{lstlisting}

  \aside{Команда идёт \lo{мимо отслеживателя}: не смешивайте в одном контексте.}
\end{frame}

% 55
\begin{frame}{AOT}
  \slidetop
  \versus{Известно при сборке}{%
    \begin{cbul}{accent}
      \item модель: компилируется заранее
      \item старт быстрее в разы
      \item запрос, написанный в одном месте
      \item предкомпиляция --- экспериментальная
    \end{cbul}}{Только при выполнении}{%
    \begin{cbul}{alt}
      \item запрос из необязательных фильтров
      \item условия внутри выражений
      \item вызовы методов в условиях
      \item перевод дерева выражений
    \end{cbul}}

  \slidegap
  \hrulethin{0.42}
  \rulegap
  {\large Пишите запросы \hi{просто}: дайте провайдеру перевести.}
\end{frame}

% 56
\begin{frame}[plain]
  \statement{ORM превращает объекты в \lo{мешки данных}, а поведение уезжает в
    сервисы.\par\rulegap
    {\normalsize\color{muted}Егор Бугаенко, «ORM Is an Offensive Anti-Pattern»}}
\end{frame}

% 57
\begin{frame}{Что с этим делать}
  \slidetop
  \versus{Чтение}{%
    \begin{cbul}{accent}
      \item проекция в DTO
      \item без сущностей вовсе
      \item без отслеживания
    \end{cbul}}{Запись}{%
    \begin{cbul}{alt}
      \item загрузить агрегат целиком
      \item вызвать метод с правилом
      \item сохранить рядом
    \end{cbul}}

  \slidegap
  \hrulethin{0.42}
  \rulegap
  {\large ORM --- генератор SQL и маппинга, \hi{а не архитектура}.}

  \rulegap
  \aside{Агрегат --- то, что меняется вместе: заказ с позициями. Границы агрегата
    --- границы транзакции.}
\end{frame}

% 58
\begin{frame}{Альтернативы}
  \slidetop
  {\renewcommand{\arraystretch}{1.5}
  \begin{tabular}{@{}L{3.0cm}L{4.9cm}L{4.9cm}@{}}
    & \colhead{accent}{что это} & \colhead{alt}{когда} \\
    \hi{\textbf{EF Core}} & LINQ, отслеживание, миграции      & запросов мало, и они
      простые \\
    \textbf{Dapper}       & вы пишете SQL, он делает маппинг  & отчёты, аналитика,
      оконные функции \\
    \textbf{linq2db}      & тот же LINQ, ближе к SQL, без отслеживания & команда живёт
      в SQL и не хочет магии \\
  \end{tabular}}

  \slidegap
  \aside{Часто вместе: EF Core --- запись и простое чтение, Dapper --- сложное.
    Все работают через \hi{один драйвер}.}
\end{frame}

% 59
\begin{frame}[plain]
  \keyline{Меняйте и сохраняйте рядом, массовое --- одной командой, запросы
    пишите просто.\\ ORM --- генератор, а не архитектура.}
\end{frame}

\end{document}
